%%%%%%%%%%%%%%%%%%%%%
% Documento maestro %
%%%%%%%%%%%%%%%%%%%%%
\documentclass{uanl}

%%%%%%%%%%%%%%%%%%%%%%%%%%%%%%%%%%%%%%%%%%%
% Cargando paquetes y definiendo opciones %
%%%%%%%%%%%%%%%%%%%%%%%%%%%%%%%%%%%%%%%%%%%
% Aquí puedes cargar los paquetes que vas a usar. La clase
% fime ya incluye babel, inputenc, graphicx y los de la AMS.
% Cargar un paquete está a tu libertad (y responsabilidad).
\usepackage{hyperref}
    \hypersetup{breaklinks=true,colorlinks=true,linkcolor=black,citecolor=black,urlcolor=black}

% Ajusta la altura del número de página
\setlength{\footskip}{0.5cm}

%%%%%%%%%%%%%%%%%%%%%
% Definiendo campos %
%%%%%%%%%%%%%%%%%%%%%
\def\titulo{Título de la tesis}
\def\autor{Juan Jesús Torres Solano}
\def\matricula{2173262}
\def\grado{Maestría en Ciencia de Datos}
% En caso de que el grado tenga orientación o especialidad llenar el siguiente
% campo dejando un ESPACIO INICIAL, en caso contrario, dejar vacío
%\def\orientacion{Con Orientación en Fibras Ópticas, Fotónica y Sensores Ópticos.}
% Coloca el mes con mayúscula inicial
\def\fecha{marzo 2025}

\def\asesor{Dr. José Anastasio Hernández Saldaña}
\def\revisorA{Nombre del revisor A}
\def\revisorB{Nombre del revisor B}
% En el caso de que tu tesis sea de doctorado activa la variable cambiándola a \doctoradotrue
% y define tus otros dos revisores
\newif\ifdoctorado\doctoradotrue
\def\revisorC{Nombre del revisor C}
\def\revisorD{Nombre del revisor D}
% El visto bueno siempre va
\def\vobo{Dra. Azucena Yoloxóchitl Ríos Mercado}

%%%%%%%%%%%%%%%%%%%%%%%
% Inicia el documento %
%%%%%%%%%%%%%%%%%%%%%%%
\begin{document}

\frontmatter
\pagestyle{main}

%%%%%%%%%%%%%%%%%%%%%%%%
% Primer portada: UANL %
%%%%%%%%%%%%%%%%%%%%%%%%
\thispagestyle{empty}

\begin{scshape}
\begin{center}
	{\Large \uanl} \\[5mm]
	{\large \fcfm} \\[5mm]
	{\large \mcd}
	\vskip15mm
	\includegraphics[height=55mm]{uanl}
	\vskip12mm
	\begin{tabular}{p{11cm}}
		\centering
		{\large \titulo}
	\end{tabular}
	\vskip7mm
	{por}\\[7mm]
	{\large \autor}\\[7mm]
        % Si en tu posgrado la tesis no es opcional (como sí lo es en licenciatura), no modifiques esta línea:
	{como requisito parcial para obtener el grado de}\\[3mm]
	\MakeUppercase{\grado}\\
	\orientacion
	\vfill
	\fecha\\
	\uanl\\
	\fcfm
\end{center}
\end{scshape}

%%%%%%%%%%%%%%%%%%%%%%%%%%%%%
% Carta del comité de tesis %
%%%%%%%%%%%%%%%%%%%%%%%%%%%%%
\newpage
\thispagestyle{empty}
\enlargethispage{5mm}

{\renewcommand{\baselinestretch}{1.1}\selectfont
\begin{center}\vspace*{-25mm}\hspace*{-10mm}
\begin{minipage}{170.5mm}
\hspace{-1.5mm}\includegraphics[height=20mm]{uanl}\hfill\raise0mm\hbox{\includegraphics[height=20mm]{fcfm}}
\hrule\vspace{0.5mm}
\scalebox{.5}{\MakeUppercase{\uanl}}\hfill\scalebox{.5}{\MakeUppercase{\fcfm}}\medskip
\end{minipage}
%\vskip4mm{\sc\large\uanl\\\fcfm\\[3pt]\pifi}\vskip6mm
\end{center}

Los miembros del Comité de Tesis recomendamos que la Tesis <<\titulo>>, realizada por el alumno \autor, con número de matrícula \matricula, sea aceptada para su defensa como requisito parcial para obtener el grado de \grado\orientacion.
\ifdoctorado\vskip10mm\else\vskip8mm\fi

\begin{center}
El Comité de Tesis\\
\ifdoctorado\vskip15mm\else\vskip25mm\fi

\ifdoctorado{%%%
\begin{tabular}{p{37mm}p{21mm}p{12mm}p{21mm}p{37mm}}
	\cline{2-4}
	& \multicolumn{3}{c}{\asesor} & \\
	& \multicolumn{3}{c}{Asesor}  & \\[15mm]
	\cline{1-2} \cline{4-5}
	%\multicolumn{2}{c}{\revisorA} & & \multicolumn{2}{c}{\revisorB} \\
	%\multicolumn{2}{c}{Revisor}   & & \multicolumn{2}{c}{Revisor}   \\[17mm]
	\cline{1-2} \cline{4-5}
	\multicolumn{2}{c}{\revisorC} & & \multicolumn{2}{c}{\revisorD} \\
	\multicolumn{2}{c}{Revisor}   & & \multicolumn{2}{c}{Revisor}   \\[2mm]
	& \multicolumn{3}{c}{Vo. Bo.} & \\[14mm]
	\cline{2-4}
	& \multicolumn{3}{c}{\vobo}   & \\
	& \multicolumn{3}{c}{Subdirector de Estudios de Posgrado}   & \\ &&&&
\end{tabular}
}\else{%%%
\begin{tabular}{p{37mm}p{21mm}p{12mm}p{21mm}p{37mm}}
	\cline{2-4}
	& \multicolumn{3}{c}{\asesor} & \\
	& \multicolumn{3}{c}{Asesor}  & \\[19mm]
	\cline{1-2} \cline{4-5}
	\multicolumn{2}{c}{\revisorA} & & \multicolumn{2}{c}{\revisorB} \\
	\multicolumn{2}{c}{Revisor}   & & \multicolumn{2}{c}{Revisor}   \\[2mm]
	& \multicolumn{3}{c}{Vo. Bo.} & \\[17mm]
	\cline{2-4}
	& \multicolumn{3}{c}{\vobo}   & \\
	& \multicolumn{3}{c}{Coordinadora de la Maestría en Ciencia de Datos}   & \\ &&&&
\end{tabular}
}\fi%%%

\vfill

\snnl, \MakeLowercase{\fecha}

\end{center}}

% Dedicatoria

\thispagestyle{empty}
\vspace*{17mm}

\begin{flushright}
\begin{itshape}

Aquí puedes poner tu dedicatoria\\
si es que tienes una.\bigskip\bigskip

\end{itshape}
\end{flushright}



\tableofcontents
\listoffigures
\listoftables

%Agradecimientos

\chapter{Agradecimientos}
\markboth{Agradecimientos}{}

Aquí puedes poner tus agradecimientos. (No olvides agradecer a tu comité de tesis, a tus profesores, a la facultad y a CONACyT en caso de que hallas sido beneficiado con una beca).

%Resumen

\chapter{Resumen}
\markboth{Resumen}{}

{\renewcommand{\baselinestretch}{1.1}\selectfont
{\setlength{\leftskip}{10mm}
\setlength{\parindent}{-10mm}

\autor.

Candidato para obtener el grado de \grado\orientacion.

\uanl.

\fcfm.

Título del estudio: \textsc{\titulo}.

\noindent Número de páginas: \pageref*{lastpage}.}

%%% Comienza a llenar aquí
\paragraph{Objetivos y método de estudio:}
Aquí debes poner tus objetivos y métodos de estudio. (Este es el formato).

\paragraph{Contribuciones y conlusiones:}
Y aquí tus contribuciones y conclusiones. (También es parte del formato).

\bigskip\noindent\begin{tabular}{lc}
\vspace*{-2mm}\hspace*{-2mm}Firma del asesor: & \\
\cline{2-2} & \hspace*{1em}\asesor\hspace*{1em}
\end{tabular}}



\mainmatter
\pagestyle{uanl}

%%% Haz un documento para cada capítulo
\chapter{INTRODUCCIÓN}
Durante las etapas iniciles de proyectos de ingenieria civil, vial, petrolera, geohidrologica, se requiere conocer el medio sobre el cual se desarrollaran los proyectos, ya sea para establecer medidad de accion de excabacion, aprobechamiento o explotacion de algun mineral, para ello se requiere trabajar 
que es la geofisica:

La exploración geofísica consiste en un conjunto de metodologías que a través de la medición de propiedades petrofisicas del subsuelo es 

eL geoelectrica

para que sirve


oportunidades de la cuiencia de datos en el análisis en tiempo real de la respuesta 




\chapter{DELIMITACIÓN Y PLANTEAMIENTO DEL PROBLEMA}

Durante un trabajo de prospección geofísica, al realizar adquisición de datos geoeléctricos \textit{in situ}, no es posible conocer el resultado del trabajo hasta una vez realizado procesos de inversión de datos resultado de prospección, por lo que no se tiene certeza de si una lectura presenta inconsistencias o si forma parte de una respuesta esperada para determinado resultado.\\

La problemática radica en que el muestreo planeado, en la etapa de planeación de adquisición, puede no ser efectivo o no reflejar una distribución esperada, debido a que el medio que se prospecta carece de homogeneidad, siendo solo evidente durante la exploración directa del medio, lo que en muchos casos no se logra a identificar, sumando ambigüedad al proceso de interpretación.\\ 

Como primera aproximación a una solución se considerara un primer supuesto, en el cual se cuenta con respuestas de sondeos eléctricos verticales, más simple pero no menos complejo de interpretar, al contar con un solo segmento de señal, dentro de esta señal podemos encontrar distintas unidades geológicas, profundidad de acuífero, espesores de unidades y profundidad de exploración.\\

La propuesta de solución es mediante Aprendizaje automático (ML, por sus siglas en inglés), empleando una técnica de aprendizaje que permita identificar patrones que estén asociados a la respuesta de un modelo de inversión, de esta manera obtener un modelo pronóstico de la inversión, sirviendo de guía para incrementar la densidad de muestreo para mejorar el modelo y respuestas.

\chapter{JUSTIFICACIÓN}
Existen múltiples aplicaciones de ML y DL en el procesamiento, modelado e interpretación geofísica \citep{li2024, liu2020,el2001,wrona2018}, algunas de estas aplicaciones corresponden a implementaciones académicas, al igual es posible identificar aplicaciones comerciales implementadas en software principalmente de exploración sísmica de hidrocarburos \citep{diaferia2024high, panebianco2024automated}.\\


%Es de esperar que al implementar ciencias de datos en diversas áreas de la ciencia, se optimicen aspectos básicos, mejorando la analítica e interpretación de resultados, o bien, modelando predicciones de tendencias.

La aplicación de herramientas de ML durante la ejecutar de muestreo de Sondeo Eléctrico Vertical (SEV), en particular durante el proceso de adquisición \textit{in situ}, presenta la posibilidad de evaluar y ampliar el intervalo de muestreo original, mejorando la adquisición tradicional, permitiendo evaluar el muestreo contra la regresiones y establecer nuevos intervalos de oportunidad.\\

Durante el muestreo tradicional los intervalos se diseñan considerando el objetivo de exploración (idealmente), este análisis previo es un factor determinante, en esta etapa se establece el intervalo de muestreo que permitiría identificar el objeto de exploración, este objeto de exploración pueden ser unidades geológicas, acuíferos, fallas, zonas de fracturas, estructuras antropogénicas, etc.\\

De manera general se busca mantener un intervalo de muestreo menor a la frecuencia de ocurrencia del objetivo de estudio, por lo que el éxito de la exploración dependerá en su totalidad de la planeación previa de la adquisición, lo que implicaría conocer previamente la conformación, distribución y espesor de cada unidad, lo que resulta irrisorio.\\

De manera que un modelo entrenado permita generar múltiples regresiones e identificar regiones de interés no cubiertas por el muestreos y de esta manera generar puntos adicionales \textit{in situ} optimizando la adquisición de datos e impactando en la calidad del muestreo , mejorando el acotamiento de las frecuencias identificadas.\\

Pese a existir aplicaciones de ML implementadas en geofísica, no se identifica alguna enfocada este esta problemática en particular, sin embargo hay aplicaciones en otros campos de estudio con un enfoque en el muestreo y clasificación de datos no paramétricos \citep{entezami2022non, bkassiny2013multidimensional, shi2021non} empleando técnicas como Dirichlet Process Mixture Model (DPMM), radial basis function network (RBFN),  Multiple Point Statistics (MPS) y Bayesian Compressive Sensing (BCS).\\


porque random forest y su copmaracion con BCS

\chapter{Conclusión}

\section{Estado actual}

\section{Perspectiva a futuro}
\chapter{MARCO TEÓRICO}

\section{Método Geoléctrico}

\subsection{Principios Básicos}

\section{Sondeo Eléctrico Vertical}

\subsection{Procesado de Sondeo}

\subsection{Interpretación Hidrogeológica}

\section{Árboles de Decisión}

\section{Random Forest}

\subsection{Número de Árboles}

\subsubsection{Max Tree Depth}

\subsection{Bootstrap}

\subsection{Overfitting y Underfitting}

\subsection{OOB Score}

\section{Validación Cruzada}

\section{Métricas de Evaluación}
\chapter{METODOLOGÍA}

\section{ideas y apuntes}

faltan las referencias....

Esta investigación propone la implementación de un modelo de Random Forests (Bosques Aleatorios) para mejorar la calidad y optimización del proceso de adquisición de datos durante un Sondeo Eléctrico Vertical (SEV). El objetivo es utilizar técnicas de aprendizaje automático para clasificar las lecturas obtenidas y generar una regresión que permita sugerir intervalos de muestreo adicionales, mejorando así la exploración geofísica.

El flujo de trabajo propuesto consta de las siguientes etapas:

Adquisición de datos y preprocesamiento: En primer lugar, se realiza la adquisición de datos in situ utilizando el método de Sondeo Eléctrico Vertical (SEV), el cual consiste en medir la resistividad del terreno a diferentes profundidades con un intervalo de muestreo predefinido. Este intervalo se determina de acuerdo con el objetivo de exploración, como la identificación de unidades geológicas, acuíferos, fallas, fracturas o estructuras antropogénicas. Para garantizar que los datos sean adecuados para el análisis, se lleva a cabo un preprocesamiento de los datos, que incluye la limpieza de valores atípicos, normalización de las lecturas y manejo de valores faltantes.

Clasificación de las lecturas con Random Forests: Se emplea el algoritmo Random Forests, una técnica de aprendizaje automático no paramétrica que consiste en crear múltiples árboles de decisión que luego se combinan para mejorar la precisión del modelo. El modelo se entrena utilizando las lecturas de resistividad obtenidas de las distintas profundidades y las características asociadas (como el tipo de terreno, las propiedades geológicas conocidas o las variables del sondeo). Los árboles se entrenan para clasificar los datos en diferentes categorías, tales como las unidades geológicas presentes, los acuíferos, las fracturas, entre otros. La clasificación permitirá identificar patrones en las lecturas de resistividad que corresponden a diferentes tipos de formaciones geológicas.

Generación de la regresión para optimizar el intervalo de muestreo: Una vez que se haya realizado la clasificación, se utiliza el modelo de regresión basado en Random Forests para predecir la resistividad eléctrica en profundidades no muestreadas. Esto permitirá estimar la resistividad del terreno en puntos específicos que no han sido cubiertos por el muestreo inicial. A partir de estas predicciones, se podrán proponer intervalos de muestreo adicionales in situ que mejoren la representación de las formaciones geológicas de interés. La regresión también proporcionará un modelo predictivo que puede ajustarse dinámicamente para adaptar los intervalos de muestreo según las características del terreno y los objetivos de exploración.

Evaluación del modelo y ajuste de parámetros: Se realiza una evaluación exhaustiva del modelo mediante técnicas de validación cruzada para asegurarse de que el modelo esté bien entrenado y sea capaz de generalizar correctamente a nuevos datos. Además, se compara el rendimiento del modelo de Random Forests con otras técnicas de clasificación y regresión para determinar su eficacia en comparación con otros enfoques. Se analizan métricas como la precisión en la clasificación, el error cuadrático medio (RMSE) en la regresión y la capacidad de predicción en términos de muestreos adicionales.

Optimización y mejora continua: Finalmente, se optimizan los parámetros del modelo (como el número de árboles y la profundidad de los mismos) para mejorar la precisión y eficiencia del modelo. A medida que se incorporan nuevos datos de exploración y se obtienen más lecturas de resistividad, el modelo puede ser recalibrado y ajustado para mantener su efectividad en la identificación de objetivos geofísicos y en la optimización del intervalo de muestreo.


En la implementación de técnicas avanzadas como Machine Learning (ML), modelos como Random Forests pueden analizar patrones en los datos y predecir áreas con alta variabilidad de resistividad. Esto permite:

Ajustar dinámicamente el intervalo de muestreo durante la adquisición \textit{in situ}.
Generar muestreos adicionales en áreas críticas para aumentar la precisión.
Con este enfoque, se asegura que el intervalo de muestreo esté alineado con el teorema de Nyquist, optimizando la calidad de los datos y reduciendo la redundancia.

El uso de Random Forests permite abordar la complejidad y la variabilidad inherentes al proceso de adquisición de datos en el contexto de exploración geofísica. Este enfoque optimiza el proceso de muestreo, mejora la calidad de los datos obtenidos y proporciona una base para generar predicciones más precisas, contribuyendo así al éxito de las campañas de exploración.

\section{Variables de Entrada}
	Al observar las ecuaciones 5.1, 5.3 y 5.4, es posible identificar las variables involucradas en el calculo de la resistividad aparente, estos valores son medidos por un equipo automático o bien, de forma manual a través de la lectura directa en un resistivímetro, para lo cual se requiere de comprobaciones durante la adquisición.
		
	Como observamos en la tabla \ref{tab:T1}, se integran datos generados durante la planeación y valores medidos en la etapa de adquisición; Z= profundidad aparente de exploración; K= factor geométrico; AB/2= apertura total de muestreo entre dos;  MN= distancia entre electrodos de potencial; Pn= potencial natural; Pi= potencial inducido; I= corriente Inducida; PPn= promedio del potencial natural; PPi= promedio de potencial inducido; U= diferencia entre PPn y PPi; PI= promedio de corriente inducida; Rha= resistividad aparente ponderada.
	
	
	
% Please add the following required packages to your document preamble:
% \usepackage{graphicx}
% \usepackage[table,xcdraw]{xcolor}
% Beamer presentation requires \usepackage{colortbl} instead of \usepackage[table,xcdraw]{xcolor}
\begin{table}[]
	\resizebox{\columnwidth}{!}{%
		\begin{tabular}{c|c|c|c|c|c|l|l|l|l|l|l|l|l|l|l|l|l|l|l}
			\rowcolor[HTML]{C0C0C0} 
			\textbf{Z} & \textbf{K} & \textbf{AB/2} & \textit{\textbf{AB/5  \textgreater{}}} & \textit{\textbf{MN}} & \textit{\textbf{\textgreater AB/20}} & \multicolumn{1}{c|}{\cellcolor[HTML]{C0C0C0}\textbf{Pn}} & \multicolumn{1}{c|}{\cellcolor[HTML]{C0C0C0}\textbf{Pi}} & \multicolumn{1}{c|}{\cellcolor[HTML]{C0C0C0}\textbf{I}} & \multicolumn{1}{c|}{\cellcolor[HTML]{C0C0C0}\textbf{Pn}} & \multicolumn{1}{c|}{\cellcolor[HTML]{C0C0C0}\textbf{Pi}} & \multicolumn{1}{c|}{\cellcolor[HTML]{C0C0C0}\textbf{I}} & \multicolumn{1}{c|}{\cellcolor[HTML]{C0C0C0}\textbf{Pn}} & \multicolumn{1}{c|}{\cellcolor[HTML]{C0C0C0}\textbf{Pi}} & \multicolumn{1}{c|}{\cellcolor[HTML]{C0C0C0}\textbf{I}} & \multicolumn{1}{c|}{\cellcolor[HTML]{C0C0C0}\textbf{PPn}} & \multicolumn{1}{c|}{\cellcolor[HTML]{C0C0C0}\textbf{PPi}} & \multicolumn{1}{c|}{\cellcolor[HTML]{C0C0C0}\textbf{U}} & \multicolumn{1}{c|}{\cellcolor[HTML]{C0C0C0}\textbf{PI}} & \multicolumn{1}{c}{\cellcolor[HTML]{C0C0C0}\textbf{Rha}} \\ \hline
			\rowcolor[HTML]{FFFFFF} 
			\textbf{1.8} & \textbf{56.549} & \textbf{3} & \textit{\textbf{1.2}} & \textit{\textbf{0.5}} & \textit{\textbf{0.3}} & 1 & 6 & 15 & 0 & 6 & 15 & 0 & 5 & 15 & 0.3 & 5.7 & 5.3 & 15.0 & 20.11 \\ \hline
			\rowcolor[HTML]{FFFFFF} 
			\textbf{2.4} & \textbf{100.531} & \textbf{4} & \textit{\textbf{1.6}} & \textit{\textbf{0.5}} & \textit{\textbf{0.4}} & 12 & 19 & 3 & 0 & 25 & 17 & 5 & 24 & 13 & 5.7 & 22.7 & 17.0 & 11.0 & 176.45 \\ \hline
			\rowcolor[HTML]{FFFFFF} 
			\textbf{3} & \textbf{157.080} & \textbf{5} & \textit{\textbf{2}} & \textit{\textbf{0.5}} & \textit{\textbf{0.5}} & 8 & 54 & 73 & 8 & 52 & 85 & 8 & 40 & 83 & 8.0 & 48.7 & 40.7 & 80.3 & 80.28 \\ \hline
			\rowcolor[HTML]{FFFFFF} 
			\textbf{3} & \textbf{31.416} & \textbf{5} & \textit{\textbf{2}} & \textit{\textbf{2.5}} & \textit{\textbf{0.5}} & 52 & 508 & 42 & 52 & 648 & 72 & 52 & 606 & 21 & 52.0 & 587.3 & 535.3 & 45.0 & 476.64 \\ \hline
			\rowcolor[HTML]{FFFFFF} 
			\textbf{4.8} & \textbf{80.425} & \textbf{8} & \textit{\textbf{3.2}} & \textit{\textbf{2.5}} & \textit{\textbf{0.8}} & 48 & 60 & 6 & 48 & 62 & 7 & 51 & 59 & 4 & 49.0 & 60.3 & 11.3 & 5.7 & 160.85 \\ \hline
			\rowcolor[HTML]{FFFFFF} 
			\textbf{6} & \textbf{125.664} & \textbf{10} & \textit{\textbf{4}} & \textit{\textbf{2.5}} & \textit{\textbf{1}} & 52 & 60 & 2 & 52 & 60 & 5 & 52 & 60 & 5 & 52.0 & 60.0 & 8.0 & 4.0 & 301.59 \\ \hline
			\rowcolor[HTML]{FFFFFF} 
			\textbf{9} & \textbf{282.743} & \textbf{15} & \textit{\textbf{6}} & \textit{\textbf{2.5}} & \textit{\textbf{1.5}} & 52 & 76 & 54 & 52 & 78 & 87 & 50 & 64 & 44 & 51.3 & 72.7 & 21.3 & 61.7 & 100.04
		\end{tabular}%
	}
	\caption{Ejemplo de atributos empleados en el calculo de la resistividad aparente.}
	\label{tab:T1}
\end{table}
	
	En términos generales, el proceso de adquisición consiste en la inducción de una corriente eléctrica a través del medio geológico, dicha intensidad de corriente es registrada, junto con el valor del potencial natural (Self-Potential) y el potencial inducido, generado por la inyección de corriente a tierra, obteniendo así los elementos necesarios para calcular el valor de la resistividad, habiendo previamente planeado la configuración geométrica del arreglo.

	\subsection{Datos de entrada}
		
		Para los datos de entrada se emplearon levantamientos de SEV previamente procesados e interpretados en correlación con muestreo directo por sondeo de penetración estándar (SPT por sus siglas en inglés), correspondientes a ambiente de deposito sedimentario y flujos volcánico, en ambos caso subyaciendo a unidades sedimentarias recientes.
		
		A partir de estos resultados etiquetados, valores de resistividad aparente, es como de modela variaciones, modificando el espesor de las unidades, ya que cada muestreo de resistividad aparente integra la respuesta conjunta de las unidades que la preceden, es decir las capas geoeléctricas por arriba de la profundidad aparente de exploración, para iguar las condiciones en los modelos que integran los datos de entrenamiento, se emplea la Librería PyGIMLI, la cual esta preparada para realizar esta tarea.
		
		\subsubsection{Limpieza de datos}
		
		\subsubsection{Generación de valores SEV’s}
		
		\subsubsection{Resistividad aparente}
		
		\subsubsection{Atributos cualitativos asociados a la curva de resistividad}


	\subsection{Clasificación, transformación y escalado de los datos}
		\subsubsection{Normalización o estandarización de resistividades si se observan grandes variaciones}
		\subsubsection{Transformación logarítmica de resistividad para reducir el sesgo de valores extremos}
		\subsubsection{Codificación de categorías litológicas si se incluyen como variable adicional}

\section{Diseño del Modelo Random Forest}
	\subsection{Regresión}
	\subsection{Configuración inicial del modelo, numero de árboles, profundidad, muestras y muestra mínima}
	\subsection{Configuración y optimización de hiperparametros}

\section{Preparación del dataset para la implementación del modelo}

\section{Implementación del modelo}
\subsection{Entrenamiento del modelo}
\subsection{Mapas de probabilidad y entrenamiento de regresión}

\section{Evaluación del modelo}
	\subsection{Regresión y validación cruzada}
	\subsection{Análisis de incertidumbre}


\section{Reporte estadístico}
\chapter{RESULTADOS Y CONCLUSIONES}

\appendix
%%% Haz un documento para cada apéndice, si es que tienes
\chapter{Apéndice I}






\backmatter
\pagestyle{main}

% Aquí va la bibliografía, puedes usar el entorno de LaTeX (thebibliography)
% o la herramienta BibTeX. En caso de que optes por BibTeX, puedes usar
% alguno de los archivos de estilo (mighelbib.bst o mighelnat.bst) incluidos
% en el paquete, cuyos diseños armonizan con el diseño de tesis provisto por
% fime.cls. Para muestra, basta un botón:
\bibliographystyle{mighelnat}
\bibliography{MiBiblio}

\label{lastpage}
%%Autobiografia

\chapter*{Resumen autobiográfico}
\thispagestyle{empty}

\begin{center}
\autor

Candidato para obtener el grado de\\
\grado\\
\orientacion\bigskip

\uanl\\
\fcfm\bigskip

Tesis:\\
\textsc{\large\titulo}
\end{center}\bigskip

%Aquí va tu historia
Aquí va tu historia. Recuerda que debe incluir: lugar y fecha de nacimiento, nombre de los padres, escuelas y universidades en las que se graduó después de la preparatoria, títulos o grados obtenidos (no incluir los estudios que se están concluyendo), experiencia profesional y organizaciones profesionales a las que pertenece (no incluir lista de publicaciones).


\end{document}

